% Cover letter for the Scientometrics submission of Paper D. Numbers match content/abstract.tex and FINDINGS.md.
\documentclass[10pt]{letter}
\usepackage[margin=1.6cm,top=1.2cm,bottom=1.4cm]{geometry}
\usepackage{xurl}
\sloppy
\usepackage[hidelinks]{hyperref}
\signature{Hossein Mohasel Arjomandi\\Microsoft, Redmond, WA, USA\\\href{mailto:hosseinm@microsoft.com}{hosseinm@microsoft.com}}
\address{}
\date{\today}
\begin{document}
\begin{letter}{Editorial Office\\\emph{Scientometrics}}
\opening{Dear Editor,}

I am pleased to submit the manuscript \emph{``Do links to retracted papers predict future retractions? An
out-of-sample test on the open biomedical citation graph''} for consideration as a research article in
\emph{Scientometrics}.

\textbf{Context.} Research-integrity screening tools flag papers that cite retracted work, and a recent retrospective
study of Web of Science data found that such papers are retracted more often. That association has not been tested
the way a screen is used: with information frozen at one date and outcomes observed afterwards, and separately from
the retraction waves that sweep through individual journals.

\textbf{What the manuscript does.} Using only open data, I computed scores for 10.5 million biomedical papers on the
January 2024 OpenCitations graph with the retractions known at that date, froze them with a hash, and compared them
with retractions recorded by Retraction Watch from February 2024 to October 2026. The scores, metrics and decision
rule were specified in advance. The main findings are:
\begin{itemize}
\item Citing a paper that had already been retracted predicts later retraction: within journal and publication year
  the odds ratio is 2.99 (95\% CI 2.73--3.27), unchanged when reference-list length is also held fixed.
\item As a screen the signal is weak: 0.9\% of the 1,000 highest-ranked papers were retracted later (2.2\% once
  reviews, which often cite retracted work but are rarely retracted, are set aside), and 84\% of later retractions
  cite no retracted paper.
\item Timing matters: inside the late-2023 retraction wave, a journal's own retraction rate was the far better screen.
  Supporting analyses show that retracted paper-mill papers cluster in citation communities beyond their journal and
  year, and that NIH-funded papers cite already-retracted work as often as other papers in the same journal and year.
\end{itemize}

\textbf{Why \emph{Scientometrics}.} The work combines citation-network analysis, open bibliographic data and research
integrity, all central to the journal. It builds on studies of retracted papers and their citations published in
\emph{Scientometrics}, and it gives editors and tool builders a concrete, quantified answer to how much links to
retracted work are worth as a screening signal. Retracted papers are described only by their public labels; no
individuals are named, and no score of any non-retracted paper is released.

\textbf{Openness.} All inputs are open. The code, the specification file with its recorded hash, run records with
checksums, and a map from every number to the run that produced it are at
\mbox{\href{https://github.com/HosseinMA96/retracted-neighbours-open-data}{github.com/HosseinMA96/retracted-neighbours-open-data}}; the derived dataset is on Zenodo
(\mbox{\href{https://doi.org/10.5281/zenodo.23874477}{doi:10.5281/zenodo.23874477}}). Analyses added after the outcomes were seen are labelled post hoc.

\textbf{Related submission.} A separate manuscript by the same author, \emph{``Can open citation data detect
citation stacking? An audit against Clarivate's journal suppressions''}, is under consideration at
\emph{Scientometrics} (submitted 31 October 2026). It uses the same open citation snapshot but addresses a different
question (journal-level citation stacking); the two manuscripts share no results or text.

The manuscript is original, has not been published, and is not under consideration elsewhere. I am the sole author,
there is no funding to declare and I have no competing interests. The use of generative AI is disclosed in the
Methods section, in line with Springer Nature's policy. The views expressed are my own and do not necessarily reflect
those of Microsoft.

Thank you for considering this manuscript.

\closing{Sincerely,}
\end{letter}
\end{document}
